\documentclass[12pt]{article}
\usepackage[utf8]{inputenc}
\usepackage[T1]{fontenc}
\usepackage{hyperref}
\title{Pentlatch Language Resources for K-12 Curriculum Development}
\author{Compiled by Claude Sonnet 5.5 (claude-sonnet-5-5) for Andr\'as Kornai}
\date{}
\begin{document}
\maketitle

\section{Introduction}
Pentlatch (ISO 639-3 ptw; Glottocode pent1242) is a Central Salish language once spoken on Vancouver Island between Comox and Nanaimo. Sources report that the last fluent speaker died around 1940. The search results also say that Qualicum First Nation researchers declared it a living language in December 2023, and that Elder Bill Recalma began recording and teaching what he knew in 2017. Documentation is mainly Franz Boas's 1886 notes.

\section{Primary Educational Resources}
No curriculum or classroom materials were found in the search results.

\section{Linguistic and Reference Materials}
\begin{itemize}
\item Wikipedia, ``Pentlatch language'': \url{https://en.wikipedia.org/wiki/Pentlatch_language}
\item Andreatta, Recalma and Urbanczyk, paper on reconstructing the sounds behind Boas's transcriptions (ICSNL 59, 2024): \url{https://lingpapers.sites.olt.ubc.ca/files/2024/07/ICSNL59_Andreatta_Recalma_Urbanczyk_finished.pdf}
\item American Philosophical Society, K'\'omoks materials, ACLS collection (Boas holdings, including Comox and Pentlatch texts): \url{https://indigenousguide.amphilsoc.org/entry/40023}
\item Working Papers article from the University of Victoria journal site (appeared in the results; I did not read its content): \url{https://journals.uvic.ca/index.php/WPLC/article/view/21226/9592}
\item YDLI, ``The Pentlatch Language'': \url{https://ydli.org/langs/pentlatch.htm}
\end{itemize}

\section{Religious and Community Resources}
No religious texts were found. The revitalization work is led by Qualicum First Nation, per the search summaries; no community website was found.

\section{Digital and Audio Resources}
No audio exists from the period of the Boas notes. Other pages that appeared in results:
\begin{itemize}
\item \url{https://rosettapanglossia.longnow.org/index.php?title=Ptw}
\end{itemize}

\section{Curriculum Development Context}
A heritage and revitalization curriculum would have to be built around reconstruction from Boas's notes together with the community, and should be led by Qualicum First Nation and K'\'omoks-area knowledge holders. The 2024 paper shows that even the sound values are still being worked out, so materials should be developed with the community and linguists.
\end{document}